\documentclass{article}


\usepackage[T1]{fontenc}
\usepackage[short]{optidef}
\usepackage[caption=false,font=footnotesize]{subfig}
\usepackage{adjustbox}
\usepackage{algorithm}
\usepackage{algpseudocode}
\usepackage{amsfonts}
\usepackage{amsmath}
\usepackage{amssymb}
\usepackage{amsthm}
\usepackage{cite}
\usepackage{fullpage}
\usepackage{mathtools}
\usepackage{microtype}
\usepackage{multirow}
\usepackage{pdfpages}
\usepackage{pgfplots}
\usepackage{siunitx}
\usepackage{xr-hyper}
\usepackage{hyperref}

\externaldocument[M-]{submit}

\DeclareSIUnit{\belm}{Bm}
\DeclareSIUnit{\dBm}{\deci\belm}
\DeclareSIUnit{\beli}{Bi}
\DeclareSIUnit{\dBi}{\deci\beli}

\newcounter{reviewer}
\setcounter{reviewer}{0}
\newcounter{point}[reviewer]
\setcounter{point}{0}
\newcounter{response}[reviewer]
\setcounter{response}{0}

\let\svbibcite\bibcite
\def\bibcite#1#2{\svbibcite{#1}{R#2}}
\makeatletter
\let\svbiblabel\@biblabel
\def\@biblabel#1{\svbiblabel{R#1}}
\makeatother

\renewcommand{\theequation}
	{E\arabic{equation}}

\renewcommand{\thefigure}
	{F\arabic{figure}}

\renewcommand{\thetable}
	{T\arabic{table}}

\renewcommand{\thealgorithm}
	{A\arabic{algorithm}}

\newcommand{\reviewer}
	{\stepcounter{reviewer} \bigskip \hrule \section*{Reviewer \thereviewer}}

\renewcommand{\thepoint}
	{\thereviewer.\arabic{point}}

\renewcommand{\theresponse}
	{\thereviewer.\arabic{response}}

\newenvironment{point}
	{\refstepcounter{point} \bigskip \noindent {\textbf{Comment~\thepoint} } ---\ \itshape}
	{\par}

\newenvironment{response}
	{\refstepcounter{response} \medskip \noindent \textbf{Response:}\ }
	{\medskip}


\begin{document}
	\includepdf{letter.pdf}

	\begin{reviewer}
		\begin{point}
			This paper assumed that the perfect channel state information (CSI) of the whole system is available at the base station. However, since the IRS is in general not equipped with a radio frequency chain, the accurate CSI of the reflecting links established by the IRS is very challenging to obtain. The reviewer wonders if the proposed algorithm can also be applied to a case where the imperfect CSI of the network is available at the base station.
			\label{pt:1.1}
		\end{point}

		\begin{response}
			Thank you for pointing out the CSI acquisition issue regarding the cascaded link. The proposed passive beamforming algorithms rely on such cascaded CSIT, and recent researches proposed element-wise on/off switching \cite{M-Nadeem2019}, training sequence and reflection pattern design \cite{You2019,Kang2020}, and compressed sensing \cite{Wang2020} techniques to solve this issue. Those references have been cited in the literature review. To address the comment, we have revised the simulation results and applied the proposed algorithms to an imperfect CSIT model where the estimation of the cascaded link at subband $n$ is
			\begin{equation}
				\hat{\boldsymbol{V}}_{n} = \boldsymbol{V}_{n} + \tilde{\boldsymbol{V}}_{n}
			\end{equation}
			where $\tilde{\boldsymbol{V}}_{n}$ is the estimation error with entries following i.i.d. CSCG distribution $\mathcal{CN}(0, \epsilon_{n}^2)$. Simulation result in Fig.~\ref{fi:re_csi} demonstrates the robustness of the proposed passive beamforming algorithm to cascaded CSIT inaccuracy for broadband SWIPT with different number of IRS elements.

			\begin{figure}[!h]
				\centering
				\resizebox{0.45\columnwidth}{!}{
					% This file was created by matlab2tikz.
%
%The latest updates can be retrieved from
%  http://www.mathworks.com/matlabcentral/fileexchange/22022-matlab2tikz-matlab2tikz
%where you can also make suggestions and rate matlab2tikz.
%
\definecolor{mycolor1}{rgb}{0.00000,0.44706,0.74118}%
\definecolor{mycolor2}{rgb}{0.85098,0.32549,0.09804}%
%
\begin{tikzpicture}

\begin{axis}[%
width=4.036in,
height=3.396in,
at={(0.677in,0.458in)},
scale only axis,
xmin=0,
xlabel style={font=\color{white!15!black}},
xlabel={Per-subband rate [bps/Hz]},
ymin=0,
ylabel style={font=\color{white!15!black}},
ylabel={DC [$\mu$A]},
axis background/.style={fill=white},
xmajorgrids,
ymajorgrids,
legend style={legend cell align=left, align=left, draw=white!15!black},
title style={font=\huge}, label style={font=\huge}, ticklabel style={font=\LARGE}, legend style={font=\LARGE}
]
\addplot [color=mycolor1, line width=2.0pt, mark=o, mark options={solid, mycolor1}, forget plot]
  table[row sep=crcr]{%
6.3110063772077	0.00144291875423115\\
5.97897226639126	5.01607097240716\\
5.64680269379241	11.6755601962969\\
5.31463274810295	19.0762873352017\\
4.98246098592109	26.7499172425096\\
4.65028817548905	34.3471993727882\\
4.31812029254077	42.1675849188683\\
3.98594593559631	51.473262995708\\
3.65377161092454	61.9972082546549\\
3.32161739867344	73.7509895494681\\
2.98944776569214	86.1720241705108\\
2.65727665153938	99.1089971835335\\
2.32510930876026	112.351518130612\\
1.99296688055974	125.975938942938\\
1.6608017362893	139.681543125101\\
1.32863932374057	153.29797988612\\
0.996467946713449	167.030559756857\\
0.664311375231153	180.904443952567\\
0.332156009779947	195.553101679189\\
1.1362500483725e-06	219.520021546501\\
};
\addplot [color=mycolor1, dashed, line width=2.0pt, mark=o, mark options={solid, mycolor1}, forget plot]
  table[row sep=crcr]{%
6.30439855782817	0.00143843954000356\\
5.97221794251559	5.00373573801135\\
5.63995993828306	11.6459794359913\\
5.30779029924861	19.0230744599296\\
4.97563928229891	26.6757221147027\\
4.64356432587894	34.2582990412539\\
4.31158441798787	42.0868114923656\\
3.9795897506604	51.3204526418265\\
3.64741143238515	61.8844035220233\\
3.31521350574592	73.5242026953777\\
2.98337247296442	85.872993133194\\
2.65145183685382	98.7823396688833\\
2.31956142287505	111.987860419908\\
1.98767078598425	125.500165552539\\
1.65577956532917	139.113746211084\\
1.3239482680594	152.55240593103\\
0.992278157602162	166.129063349441\\
0.660873727630513	179.748371666217\\
0.329727591293718	194.192393908755\\
8.80076644949702e-08	217.962114160156\\
};
\addplot [color=mycolor1, dotted, line width=2.0pt, mark=o, mark options={solid, mycolor1}, forget plot]
  table[row sep=crcr]{%
5.94681722047409	0\\
5.63382728182781	3.90078401004364\\
5.32083804749954	8.97463202073686\\
5.0078502345267	14.6407939447736\\
4.69486233157755	20.5773089310291\\
4.38187225054363	26.539471728473\\
4.06888041184373	32.4360636256513\\
3.75589239315053	39.1883552588207\\
3.44290558045743	46.4701193191521\\
3.12991869591931	54.813235456183\\
2.81693352803647	63.6515034802105\\
2.5039561014384	72.7955352447124\\
2.19099162709925	82.059980887216\\
1.87801924232396	91.461098149337\\
1.56503639623363	100.823403338262\\
1.25204999991622	110.259014406231\\
0.939072487841398	119.75943393223\\
0.626081882097103	129.540665234206\\
0.313081853490835	140.147477067512\\
7.12915833101467e-05	157.894101667706\\
};
\addplot [color=mycolor2, line width=2.0pt, mark=+, mark options={solid, mycolor2}, forget plot]
  table[row sep=crcr]{%
6.74494050114811	0.0017780354320067\\
6.39007796206246	7.06012826574983\\
6.03507961899998	16.709883719072\\
5.6800814667746	27.3281970964577\\
5.32508555378385	38.1282421557938\\
4.97008943140078	48.5086971863898\\
4.61509436071278	60.8187863279741\\
4.26009528179264	75.6282966644584\\
3.90510826554925	92.3058048320567\\
3.55010840737924	110.364762698681\\
3.19510491832684	129.161018751469\\
2.84009792134078	148.778608648376\\
2.48507840671485	168.816773677303\\
2.13006206540349	189.188492238012\\
1.77506344930209	210.610260420195\\
1.42005805719262	232.238345603231\\
1.06503455587726	254.285159357222\\
0.710008190170042	275.50742470563\\
0.355009645577302	296.877358295413\\
8.34346484672978e-09	332.278155146066\\
};
\addplot [color=mycolor2, dashed, line width=2.0pt, mark=+, mark options={solid, mycolor2}, forget plot]
  table[row sep=crcr]{%
6.73514961204729	0.0017690930923067\\
6.37990167706694	7.02398417518874\\
6.02458985188032	16.6157887579869\\
5.66936750335739	27.1619469786471\\
5.31427124656903	37.8798641536591\\
4.95929565643053	48.1623563272398\\
4.60437569664923	60.4387703938643\\
4.24944262089574	75.1351291676052\\
3.89451396435008	91.7052842157096\\
3.53966452267356	109.557209476152\\
3.18488537304119	128.336543840389\\
2.83006847806064	147.711687841547\\
2.47551669022242	167.498344826428\\
2.12122319871569	187.821735860952\\
1.76663464428464	208.929770818195\\
1.41241713017586	230.13636207396\\
1.05834463617405	251.759403760308\\
0.70441573871233	272.421585140811\\
0.351122537332296	293.533549576334\\
8.74671235533785e-09	328.479178159679\\
};
\addplot [color=mycolor2, dotted, line width=2.0pt, mark=+, mark options={solid, mycolor2}, forget plot]
  table[row sep=crcr]{%
6.06568180785804	0\\
5.74643801723957	4.29465295156347\\
5.42718999725819	9.93372833440709\\
5.10794592574478	16.1994494771073\\
4.78870155627181	22.6996241943057\\
4.46945542538136	29.0291226335639\\
4.15020837437215	35.9023833041809\\
3.83096285290911	43.7977255791796\\
3.51171901611387	52.7130256331732\\
3.19247726723199	62.3354206363883\\
2.87323940756278	72.4950502974977\\
2.55400449809425	82.8039689471512\\
2.23476749764932	93.262196590131\\
1.91553024355807	103.758768614458\\
1.59629859821389	114.241917448414\\
1.27706598767081	124.478800925889\\
0.957833521031375	135.10780331253\\
0.638591366630782	145.941950305931\\
0.319334346188265	157.617727171717\\
7.47316309940721e-05	177.484025981736\\
};
\addplot [color=black, line width=2.0pt]
  table[row sep=crcr]{%
-1	-1\\
};
\addlegendentry{$\epsilon_n^2 = 0.1 \Lambda_{\mathrm{I}}\Lambda_{\mathrm{R}}$}

\addplot [color=black, dashed, line width=2.0pt]
  table[row sep=crcr]{%
-1	-1\\
};
\addlegendentry{$\epsilon_n^2 = 10 \Lambda_{\mathrm{I}}\Lambda_{\mathrm{R}}$}

\addplot [color=black, dotted, line width=2.0pt]
  table[row sep=crcr]{%
-1	-1\\
};
\addlegendentry{$\epsilon_n^2 \mathrel{\to} \infty$}

\addplot [color=mycolor1, line width=2.0pt, only marks, mark=o, mark options={solid, mycolor1}]
  table[row sep=crcr]{%
-1	-1\\
};
\addlegendentry{$L = 20$}

\addplot [color=mycolor2, line width=2.0pt, only marks, mark=+, mark options={solid, mycolor2}]
  table[row sep=crcr]{%
-1	-1\\
};
\addlegendentry{$L = 40$}

\end{axis}
\end{tikzpicture}%

				}
				\caption{Average R-E region with imperfect cascaded CSIT for $M=1$, $N=16$, $\sigma_n^2=\SI{-40}{\dBm}$, $B=\SI{10}{\MHz}$ and $d_{\mathrm{H}}=d_{\mathrm{V}}=\SI{2}{\meter}$.}
				\label{fi:re_csi}
			\end{figure}

			\label{re:1.1}
		\end{response}

		\begin{point}
			This paper considered a relatively simple system model where there is only one user in the system. However, in practice, there can be many users co-existing in the SWIPT systems and there may also be two different and independent quality of service requirements: energy harvesting requirement and information decoding requirement. It could be better if the authors can clarify if the proposed algorithm can be employed in a more general case where the base station and the IRS cooperate to serve multiple users in the same time slot.
			\label{pt:1.2}
		\end{point}

		\begin{response}
			Multi-user IRS-aided SWIPT is very interesting and should be investigated in the future, but is left for future work for the ``don't run before we can walk'' reasons explained below:

			\emph{First}, it is important to remind the reviewer that the problem of joint waveform and beamforming design for a multi-user SWIPT system without IRS has never been studied in the literature. In other words, it would be premature at this stage to address this multiuser IRS-aided SWIPT since the underlying building block of multiuser SWIPT is not known. Only once the problem of multiuser joint waveform and beamforming for multi-user SWIPT has been addressed, can we investigate the multiuser joint waveform and beamforming design for IRS-aided SWIPT. Those problems could be addressed in future research, but are not the scope of this paper.

			\emph{Second}, the emphasis is on single-user because we believe a proper understanding of the single-user case is crucial before jumping into multi-user scenarios. Our modeling of SWIPT is not "simple": not only we transmit both information and power, we deal with joint space and frequency optimization through a joint beamforming and waveform problem that is challenging due to the nonlinearity of the rectifier that induces coupling among frequencies components. This completely contrasts with any existing IRS-aided SWIPT work as discussed in the introduction \cite{M-Wu2020b,M-Tang2019,M-Wu2020a}. To reflect this aspect, we have revised the title to ``IRS-Aided SWIPT: Joint Waveform, Active and Passive Beamforming Design Under Nonlinear Harvester Model''.

			\emph{Third}, multi-user WPT leads to its own set of challenges compared with single-user WPT. Taking an non-IRS-aided WPT, frequency and spatial domains can be decoupled in single-user case, while they cannot be decoupled in the multi-user case, as shown in \cite{M-Clerckx2016a,M-Huang2017}. Consequently, single-user and multi-user WPT are designed differently in \cite{M-Clerckx2016a,M-Huang2017}. Though the SWIPT case has not been studied, the same set of problems is expected to occur between single user and multiuser case and would be even more challenging since communication is delivered simultaneously with power. Adding IRS into the picture would make it even more complicated.
			\label{re:1.2}
		\end{response}

		\begin{point}
			It is well known that after employing semidefinite programming for handling the phase shift matrix at the IRS, it is very unlikely to obtain a rank-one phase shift matrix without any further modification. The reviewer notices that the authors proposed the Gaussian randomization method to ensure a rank-one solution. Also, the convergence of the proposed overall algorithm also relies on the unit-rank solution. Therefore, to make this paper more comprehensive and convincing, it is suggested to provide more results (such as figures, tables, and data analysis) in the simulation part to show that the rank-one solution can always be obtained even without applying the Gaussian randomization. As this is a very important and interesting conclusion to the colleagues working in the same area, it could be better if the authors can further discuss, interpret, and clarify this in a remark.
			\label{pt:1.3}
		\end{point}

		\begin{response}
			Thank you for raising this issue. To prove the rank-\num{1} property of the IRS matrix $\boldsymbol{\Phi}$, we notice that for any feasible $\boldsymbol{\Phi}$ to problem~\eqref{M-op:irs}, $\mathrm{tr}(\boldsymbol{\Phi})=L+1$ always holds due to the modulus constraint~\eqref{M-co:irs_modulus}. Therefore, we add a constant term $-\mathrm{tr}(\boldsymbol{\Phi})$ to \eqref{M-ob:irs} and recast problem~\eqref{M-op:irs} as
			\begin{maxi!}
				{\boldsymbol{\Phi}}{-\mathrm{tr}(\boldsymbol{\Phi})+\tilde{z}(\boldsymbol{\Phi})}{\label{op:irs_equivalent}}{\label{ob:irs_equivalent}}
				\addConstraint{R(\boldsymbol{\Phi}) \ge \bar{R}}\label{co:irs_equivalent_rate}
				\addConstraint{\mathrm{diag}^{-1}(\boldsymbol{\Phi})=\boldsymbol{1}}\label{co:irs_equivalent_modulus}
				\addConstraint{\boldsymbol{\Phi}\succeq{\boldsymbol{0}}}\label{co:irs_equivalent_sd}
				\addConstraint{\mathrm{rank}(\boldsymbol{\Phi})=1.\label{co:irs_equivalent_rank}}
			\end{maxi!}

			By applying rank constraint relaxation, problem~\eqref{ob:irs_equivalent}--\eqref{co:irs_equivalent_sd} is convex with respect to $\boldsymbol{\Phi}$ and satisfies the Slater's condition \cite{M-Boyd2004}, thus strong duality holds. The corresponding Lagrangian function at iteration $i$ is given by
			\begin{align}\label{eq:lagrangian}
				\mathcal{L}
				&= \mathrm{tr}\left(\boldsymbol{\Phi}^{(i)}\right) - \frac{1}{2} \beta_2 \rho \mathrm{tr}\Bigl(
						(\boldsymbol{C}_{\mathrm{I},0} + \boldsymbol{C}_{\mathrm{P},0}) \boldsymbol{\Phi}^{(i)}
					\Bigr) - \frac{3}{8} \beta_4 \rho^2 \Biggl(
						4 t_{\mathrm{I},0}^{(i-1)} \mathrm{tr}\Bigl(
							\boldsymbol{C}_{\mathrm{I},0} \boldsymbol{\Phi}^{(i)}
						\Bigr) + \sum_{k=-N+1}^{N-1} 2 (t_{\mathrm{P},k}^{(i-1)})^* \mathrm{tr}\Bigl(
							\boldsymbol{C}_{\mathrm{P},k} \boldsymbol{\Phi}^{(i)}
						\Bigr)
						\nonumber\\
				&\quad+ 2 (t_{\mathrm{I},0}^{(i-1)} + t_{\mathrm{P},0}^{(i-1)}) \mathrm{tr}\Bigl(
							(\boldsymbol{C}_{\mathrm{I},0} + \boldsymbol{C}_{\mathrm{P},0}) \boldsymbol{\Phi}^{(i)}
						\Bigr) - \mathrm{tr}^2 \Bigl(
							(\boldsymbol{C}_{\mathrm{I},0} - \boldsymbol{C}_{\mathrm{P},0}) \boldsymbol{\Phi}^{(i)}
						\Bigr)
					\Biggr)\nonumber\\
				&\quad+ \mu \Biggl(
					2^{\bar{R}} - \prod_{n=1}^N \biggl(
						1 + \frac{(1-\rho) \mathrm{tr}\Bigl(
							\boldsymbol{C}_n \boldsymbol{\Phi}^{(i)}
						\Bigr)}{\sigma_n^2}
					\biggr)
				\Biggr) + \mathrm{tr}\biggl(
					\mathrm{diag}(\boldsymbol{\nu}) \odot \Bigl(
						\boldsymbol{\Phi}^{(i)} \odot \boldsymbol{I} - \boldsymbol{I}
					\Bigr)
				\biggr) - \mathrm{tr} \Bigl(
					\boldsymbol{\Upsilon} \boldsymbol{\Phi}^{(i)}
				\Bigr) + \zeta,
			\end{align}
			where $\mu$, $\boldsymbol{\nu}$, $\boldsymbol{\Upsilon}$ denote respectively the scalar, vector and matrix Lagrange multiplier associated with constraint~\eqref{co:irs_equivalent_rate}, \eqref{co:irs_equivalent_modulus} and \eqref{co:irs_equivalent_sd}, and $\zeta$ collects all terms irrelevant to $\boldsymbol{\Phi}^{(i)}$. The Karush–Kuhn–Tucker (KKT) conditions on the primal and dual solutions are
			\begin{subequations}
				\begin{equation}\label{eq:lagrange_multiplier}
					\mu^\star \ge 0, \boldsymbol{\Upsilon}^\star \succeq \boldsymbol{0},
				\end{equation}
				\begin{equation}\label{eq:complementary_slackness}
					\boldsymbol{\nu}^\star \odot \mathrm{diag}^{-1}(\boldsymbol{\Phi}^\star) = \boldsymbol{0}, \boldsymbol{\Upsilon}^\star \boldsymbol{\Phi}^\star = \boldsymbol{0},
				\end{equation}
				\begin{equation}\label{eq:gradient}
					\nabla_{\boldsymbol{\Phi}^\star} \mathcal{L} = 0.
				\end{equation}
			\end{subequations}

			We then derive the gradient explicitly and rewrite \eqref{eq:gradient} as
			\begin{equation}
				\boldsymbol{\Upsilon}^\star = \boldsymbol{I} - \boldsymbol{\Delta}^\star
			\end{equation}
			with $\boldsymbol{\Delta}^\star$ given by
			\begin{align}\label{eq:delta}
				\boldsymbol{\Delta}^\star
				&= \frac{1}{2} \beta_2 \rho (\boldsymbol{C}_{\mathrm{I},0}+\boldsymbol{C}_{\mathrm{P},0}) + \frac{3}{4} \beta_4 \rho^2
					\Biggl(
						2 t_{\mathrm{I},0}^{(i-1)} \boldsymbol{C}_{\mathrm{I},0} + \sum_{k=-N+1}^{N-1} (t_{\mathrm{P},k}^{(i-1)})^* \boldsymbol{C}_{\mathrm{P},k} + (t_{\mathrm{I},0}^{(i-1)} + t_{\mathrm{P},0}^{(i-1)}) (\boldsymbol{C}_{\mathrm{I},0} + \boldsymbol{C}_{\mathrm{P},0})\nonumber\\
				&\quad- (\boldsymbol{C}_{\mathrm{I},0} - \boldsymbol{C}_{\mathrm{P},0}) \mathrm{tr}\Bigl(
							(\boldsymbol{C}_{\mathrm{I},0} + \boldsymbol{C}_{\mathrm{P},0}) \boldsymbol{\Phi}^\star
						\Bigr)
					\Biggr) + \mu^\star \sum_{n=1}^N \frac{(1-\rho) \boldsymbol{C}_n}{\sigma_n^2} \prod_{n'=1,n' \ne n}^N \Biggl(
					1 + \frac{(1-\rho)\mathrm{tr}\Bigl(
						\boldsymbol{C}_{n'} \boldsymbol{\Phi}^\star
					\Bigr)}{\sigma_{n'}^2}
				\Biggr)\nonumber\\
				&\quad- \mathrm{diag}(\boldsymbol{\nu^\star}).
			\end{align}

			Note that \eqref{eq:complementary_slackness} suggests $\mathrm{rank}(\boldsymbol{\Upsilon}^\star)+\mathrm{rank}(\boldsymbol{\Phi}^\star) \le L+1$. By reusing the proof in Appendix A of \cite{Xu2020}, we conclude $\mathrm{rank}(\boldsymbol{\Upsilon}^\star) \ge L$. On the other hand, $\boldsymbol{\Phi}^\star$ cannot be zero matrix and $\mathrm{rank}(\boldsymbol{\Phi}^\star) \ge 1$. Therefore, any optimal solution $\boldsymbol{\Phi}^\star$ to the relaxed problem~\eqref{op:irs_equivalent} satisfies $\mathrm{rank}(\boldsymbol{\Phi}^\star)=1$. Due to the equivalence between \eqref{M-ob:irs} and \eqref{ob:irs_equivalent}, $\boldsymbol{\Phi}^\star$ is also optimal to the relaxed problem~\eqref{M-op:irs}. We have clarified this point in Proposition~\ref{M-pr:relaxation}.
			\label{re:1.3}
		\end{response}

		\begin{point}
			The figures in the current version are relatively small, it is suggested to provide larger figures to help the readers better understand the results.
			\label{pt:1.4}
		\end{point}

		\begin{response}
			We sincerely apologize for the inconvenience. The figures have been amplified in the revised manuscript.
			\label{re:1.4}
		\end{response}

		\begin{point}
			In the current version, the authors claimed that by applying semidefinite relaxation and omitting the rank-one constraint, the performance loss is negligible. The reviewer wonders if this is because of the relatively simple system model, as there is only one single user who has both power and information requirements. It could be better if the authors can further discuss this issue for a more general multiple user scenario.
			\label{pt:1.5}
		\end{point}

		\begin{response}
			We agree with the reviewer that the relaxation on the rank-one constraint deserves more attention. However, as discussed in Response~\ref{re:1.2}, multi-user IRS-aided SWIPT is out the scope of this paper and we would like to provide some general ideas to address this issue. For any passive beamforming problem with non-convex unit-rank constraint, if the proof to Proposition~\ref{M-pr:relaxation} is unavailable, existing approaches include 1) relax the constraint and use Gaussian randomization method to obtain a high-quality solution \cite{M-Luo2010}, 2) replace the unit-rank constraint with the constraint on largest singular value $\mathrm{Tr}(\boldsymbol{\Phi})-\sigma_{\max} (\boldsymbol{\Phi}) \le 0$ then solve the problem by the penalty method \cite{M-Wu2020a}, \cite{Wu2021}, 3) apply the gradient descent in the manifold space \cite{Zuo2020}. All three techniques are expected to provide a unit-rank solution with performance very close to the original solution.
			\label{re:1.5}
		\end{response}
	\end{reviewer}

	\begin{reviewer}
		\begin{point}
			New RIS models are now being adopted as in \cite{M-Abeywickrama2020} where it has been shown that the reflected signals depend on the direction of the arriving signal and this needs to be included in the analysis for realistic quantification.
			\label{pt:2.1}
		\end{point}

		\begin{response}
			Thank you for sharing this paper. It investigated the impact of non-zero effective resistance on the reflection pattern and pointed out that the amplitude of the reflection coefficient depends on the phase shift forced on the incoming signal when power dissipation is considered at the IRS. It also proposed an analytical IRS model together with an BCD algorithm to maximize the achievable rate by passive beamforming. Simulation results emphasized the importance of modeling such a relationship in practical IRS design. There exist various refined models of IRS in the literature, however, we have finally decided to use the most common and simplest IRS model at the current stage to reduce the design complexity and provide a primary benchmark for practical IRS-aided SWIPT. We have mentioned in the conclusion that ``Also, the use of more involved IRS architecture based on direction- and frequency-dependent reflection \cite{M-Abeywickrama2020} and/or partially/fully-connected reflection \cite{M-Shen2020a} could be considered in future works.''
			\label{re:2.1}
		\end{response}

		\begin{point}
			Why is MRT considered as precoder by \eqref{M-eq:b_n} rather than optimizing it? Is it globally optimal too?
			\label{pt:2.2}
		\end{point}

		\begin{response}
			In the single-user scenario, the global optimal information and power precoders coincide at MRT. To prove this, we decouple the waveform in the spatial and frequency domains by
			\begin{equation}\label{eq:w}
				\boldsymbol{w}_{\mathrm{I/P}, n} = s_{\mathrm{I/P}, n} \boldsymbol{b}_{\mathrm{I/P}, n}
			\end{equation}
			where $s_{\mathrm{I/P},n}$ denotes the amplitude of modulated/multisine waveform at tone $n$, and $\boldsymbol{b}_{\mathrm{I/P}, n}$ denotes the information/power precoder. The MRT precoder at subband $n$ is given by
			\begin{equation}\label{eq:b}
				\boldsymbol{b}_{\mathrm{I/P}, n}^\star = \frac{\boldsymbol{h}_n}{\lVert{\boldsymbol{h}_n}\rVert}
			\end{equation}

			From the perspective of WIT, the MRT precoder maximizes $\lvert{\boldsymbol{h}_{n}^H \boldsymbol{w}_{\mathrm{I}, n}}\rvert = \lVert{\boldsymbol{h}_{n}}\rVert s_{\mathrm{I}, n}$ thus maximizes the rate (8). From the perspective of WPT, the MRT precoder maximizes $(\boldsymbol{h}_{n}^H \boldsymbol{w}_{\mathrm{I/P}, n})(\boldsymbol{h}_{n}^H \boldsymbol{w}_{\mathrm{I/P}, n})^* = \lVert{\boldsymbol{h}_{n}}\rVert^2 s_{\mathrm{I/P}, n}^2$ thus maximizes the second and fourth order DC terms \eqref{M-eq:y_I2}--\eqref{M-eq:y_P4}. Hence, MRT is the global optimal active precoder and no dedicated energy beams are required in the single-user SWIPT. We have separated the active beamforming design from the waveform design in the revised manuscript, summarized the conclusion in Proposition~\ref{M-pr:mrt}, and add the proof above to clarify this point.
			\label{re:2.2}
		\end{response}

		\begin{point}
			Some strong assumptions like perfect CSI availability limit the practical utility of the proposed analytical results.
			\label{pt:2.3}
		\end{point}

		\begin{response}
			The reviewer is referred to Response \ref{re:1.1}. Indeed, the assumption of perfect CSIT is very ideal and the existing channel estimation protocols may not provide decent results in practice. We have investigated the impact of CSIT estimation error of the cascaded link on the R-E performance, and Fig.~\ref{fi:re_csi} shows that the proposed algorithms are robust to CSIT inaccuracy for different $L$.
			\label{re:2.3}
		\end{response}

		\begin{point}
			All the assumptions and relaxations adopted used in the derivation of results as in \eqref{M-eq:z_irs_approx} need to be explicitly mentioned along with appropriate justification for the same.
			\label{pt:2.4}
		\end{point}

		\begin{response}
			We appreciate your suggestion and have revised the manuscript correspondingly. The original objective function \eqref{M-eq:z_irs} is differentiable and non-concave in $\mathbb{C}^{4N - 2}$, and we approximate (linearize) the second-order terms by the first-order Taylor expansions \eqref{M-eq:taylor_1}--\eqref{M-eq:taylor_3} to formulate SCA problems of the original passive beamforming problem (i.e. maximize \eqref{M-eq:z_irs} s.t. \eqref{M-co:irs_rate}--\eqref{M-co:irs_rank}). Also, \eqref{M-eq:coupled_terms} holds due to symmetry. Hence, the quadratic objective function \eqref{M-eq:z_irs_approx} is obtained by plugging \eqref{M-eq:taylor_1}--\eqref{M-eq:coupled_terms} into \eqref{M-eq:z_irs}. Solving relaxed problem \eqref{M-op:irs} iteratively is guaranteed to converge to a local optimal point of the original passive beamforming problem, and the proof is provided in Appendix~\ref{M-ap:sca}.
			\label{re:2.4}
		\end{response}

		\begin{point}
			Some transformations have been made while solving the original problem, but it has not been explicitly mentioned whether it is equivalent to transformation or not.
			\label{pt:2.5}
		\end{point}

		\begin{response}
			Thank you for the reminder. All transformations are equivalent to their original form and we have mentioned this point explicitly in the revised manuscript.
			\label{re:2.5}
		\end{response}

		\begin{point}
			Are the proposed solutions locally optimal or globally optimal? It is not clear whether the convergence of proposed solution methodologies is local or global? Also, how fast is it?
			\label{pt:2.6}
		\end{point}

		\begin{response}
			Algorithm~\ref{M-al:sca}--\ref{M-al:bcd} provide local optimal solutions with local convergence proof, and the performance indeed depends on the initialization. For the passive beamforming Algorithm~\ref{M-al:sca} and \ref{M-al:m_sca}, we initialize the phase shift of all IRS elements as i.i.d. uniform random variables over $[0, 2\pi)$. For Algorithm~\ref{M-al:gp}, we initialize the modulated waveform by the Water-Filling (WF) strategy \ref{eq:wf} and the multisine waveform by the Scaled Matched Filter (SMF) scheme \ref{eq:smf}, and assume $\rho^{(0)}=\bar{\rho}^{(0)}=1$. These parameters are used for general initialization and regulated automatically afterwards. However, as Algorithm~\ref{M-al:sca}, \ref{M-al:gp} and \ref{M-al:m_sca} only converge to local optimal solutions, few R-E points might be strictly worse (with less rate and energy) than the others especially when $N$ is very large. To address this issue, we draw the R-E boundary from the high-rate low-energy (lower right) points to the low-rate high-energy (upper left) points. If a point is strictly dominated, we discard the candidate, reinitialize Algorithm~\ref{M-al:bcd} and \ref{M-al:lc_bcd} by the solution at the previous point, then perform the optimization again. For a tolerance of $\epsilon=10^{-8}$, Algorithm~\ref{M-al:sca}--\ref{M-al:lc_bcd} typically converge within 2, 7, 3, 2, 3 iterations, respectively.

			We also would like to point out that, even in the absence of IRS, the problem of joint waveform and beamforming design for SWIPT is also non-convex and algorithms only converging toward a local optimum have been identified, with the global optimum remaining unsolved.
			\label{re:2.6}
		\end{response}

		\begin{point}
			The time complexity of the proposed algorithms, especially involving branch and bound methods, seems to be high especially applications assuming perfectly CSI availability as the coherence times are practically pretty low. So, the authors would like to justify it so that the proposed solution can be obtained over relatively short coherence intervals.
			\label{pt:2.7}
		\end{point}

		\begin{response}
			Problem \eqref{M-op:irs} is not a Semidefinite Programming (SDP) due to the squared term in the objective function. Hence, existing complexity analysis tools cannot be direct applied to Algorithm~\ref{M-al:sca} and \ref{M-al:m_sca}. For Algorithm~\ref{M-al:gp}, the computational complexity scales exponentially with the number of subbands \cite{M-Huang2017,M-Chiang2005} but an analytical expression cannot be derived on top of relevant literatures. To facilitate practical SWIPT implementation and address the reviewer's comment, we have added in the revised manuscript two closed-form adaptive waveform amplitude schemes by combining WF and SMF in the time and power domains. For WIT, the optimal WF strategy assigns the amplitude of modulated tone $n$ by
			\begin{equation}\label{eq:wf}
				s_{\mathrm{I}, n} = \sqrt{2\left(\lambda - \frac{\sigma_n^2}{P \lVert{\boldsymbol{h}_n}\rVert^2}\right)^+}
			\end{equation}
			where $\lambda$ is chosen to satisfy the power constraint $\lVert{\boldsymbol{s}_I}\rVert^2 / 2 \le P$. The closed-form solution can be obtained by iterative power allocation \cite{M-Tse2005}, and the details are omitted here. On the other hand, SMF was proposed in \cite{M-Clerckx2017} as a suboptimal WPT resource allocation scheme that assigns the amplitude of sinewave $n$ by
			\begin{equation}\label{eq:smf}
				s_{\mathrm{P}, n} = \sqrt{\frac{2 P}{\sum_{n=1}^N \lVert{\boldsymbol{h}_n \rVert^{2 \alpha}}}}\lVert{\boldsymbol{h}_n}\rVert^\alpha
			\end{equation}
			where the scaling ratio $\alpha \ge 1$ is given and can be adjusted to exploit the rectifier nonlinearity and frequency selectivity. When the receiver works in TS mode, there is no superposition in the suboptimal waveform design. Modulated waveform with amplitude \eqref{eq:wf} is used in the data session while multisine waveform with amplitude \eqref{eq:smf} is used in the energy session. When the receiver works in PS mode, we jointly design the combining ratio $\delta$ with the splitting ratio $\rho$, and assign the superposed waveform amplitudes as
			\begin{align}
				s_{\mathrm{I}, n} &= \sqrt{2(1 - \delta)\left(\lambda - \frac{\sigma_n^2}{P \lVert{\boldsymbol{h}_n}\rVert^2}\right)^+}, \label{eq:s_i}\\
				s_{\mathrm{P}, n} &= \sqrt{\frac{2 \delta P}{\sum_{n=1}^N \lVert{\boldsymbol{h}_n \rVert^{2 \alpha}}}}\lVert{\boldsymbol{h}_n}\rVert^\alpha. \label{eq:s_p}
			\end{align}

			Besides, minor modifications are required for passive beamforming to accommodate both low-complexity waveform schemes. To achieve the WIT point (i.e., $\eta=\rho=0$), the rate \eqref{M-eq:R_irs} instead of the DC current \eqref{M-eq:z_irs_approx} should be maximized. To achieve any non-WIT point (i.e., $z \ne 0$), the rate constraint \eqref{M-co:irs_rate} should be dropped as the R-E region is instead characterized by varying either $\eta$ or $\delta$ and $\rho$. The Modified-SCA (M-SCA) Algorithm~\ref{al:m_sca} summarizes the modified passive beamforming design when the receiver works in PS mode. The tightness of rank relaxation and local optimality are proved similar to Proposition~\ref{M-pr:relaxation} and \ref{M-pr:sca} thus omitted here.
			\begin{algorithm}[!h]
				\caption{M-SCA: IRS Phase Shift.}
				\label{al:m_sca}
				\begin{algorithmic}[1]
					\State \textbf{Input} $\beta_2,\beta_4,\boldsymbol{h}_{\mathrm{D},n},\boldsymbol{V}_{n},\sigma_n,\boldsymbol{w}_{\mathrm{I/P},n},\rho,\epsilon$, $\forall n$
					\State Construct $\boldsymbol{V},\boldsymbol{M},\boldsymbol{M}_n,\boldsymbol{C}_{n},\boldsymbol{C}_{\mathrm{I/P},k}$, $\forall n,k$
					\State \textbf{Initialize} $i \gets 0$, $\boldsymbol{\Phi}^{(0)}$
					\If{$\rho=0$}
						\State Get $\boldsymbol{\Phi}^{\star}$ by maximizing \eqref{M-eq:R_irs} s.t. \eqref{M-co:irs_modulus}, \eqref{M-co:irs_sd}
					\Else
						\State Set $t_{\mathrm{I/P},k}^{(0)}$, $\forall k$ by \eqref{M-eq:t_k}
						\Repeat
							\State $i \gets i + 1$
								\State Get $\boldsymbol{\Phi}^{(i)}$ by maximizing \eqref{M-eq:z_irs_approx} s.t. \eqref{M-co:irs_modulus}, \eqref{M-co:irs_sd}
								\State Update $t_{\mathrm{I/P},k}^{(i)}$, $\forall k$ by \eqref{M-eq:t_k}
								\State Compute $z^{(i)}$ by \eqref{M-eq:z_irs}
						\Until $\lvert z^{(i)}-z^{(i-1)} \rvert \le \epsilon$
						\State Set $\boldsymbol{\Phi}^{\star}=\boldsymbol{\Phi}^{(i)}$
					\EndIf
					\State Get $\bar{\boldsymbol{\phi}}^\star$ by eigen decomposition, $\boldsymbol{\Phi}^{\star}=\bar{\boldsymbol{\phi}}^\star(\bar{\boldsymbol{\phi}}^\star)^H$
					\State Set $\theta_l^\star=\arg(\bar{\phi}_l^\star/\bar{\phi}_{L+1}^\star)$, $\forall l$, $\boldsymbol{\phi}^{\star}=[e^{j\theta_1^\star},\dots,e^{j\theta_L^\star}]^H$
					\State \textbf{Output} $\boldsymbol{\phi}^{\star}$
				\end{algorithmic}
			\end{algorithm}
			For the low-complexity design, when the receiver works in PS mode, we instead obtain the phase shift by Algorithm~\ref{al:m_sca} and the waveform amplitude by \eqref{eq:s_i} and \eqref{eq:s_p}. The steps are summarized in the Low Complexity-BCD (LC-BCD) Algorithm~\ref{al:lc_bcd}.
			\begin{algorithm}[!h]
				\caption{LC-BCD: Waveform and Beamforming.}
				\label{al:lc_bcd}
				\begin{algorithmic}[1]
					\State \textbf{Input} $\beta_2,\beta_4,\boldsymbol{h}_{\mathrm{D},n},\boldsymbol{V}_{n},P,\sigma_n,\delta,\rho,\epsilon$, $\forall n$
					\State \textbf{Initialize} $i \gets 0$, $\boldsymbol{\phi}^{(0)},\boldsymbol{b}_{\mathrm{I/P},n}^{(0)},\boldsymbol{s}_{\mathrm{I/P}}^{(0)}$, $\forall n$
					\State Set $\boldsymbol{w}_{\mathrm{I/P},n}^{(0)}$, $\forall n$ by \eqref{M-eq:w}
					\Repeat
						\State $i \gets i + 1$
						\State Get $\boldsymbol{\phi}^{(i)}$ based on $\boldsymbol{w}_{\mathrm{I/P}}^{(i-1)}$ by Algorithm~\ref{al:m_sca}
						\State Update $\boldsymbol{h}_n^{(i)},\boldsymbol{b}_n^{(i)}$, $\forall n$ by \eqref{M-eq:h_n}, \eqref{M-eq:b_n}
						\State Update $\boldsymbol{s}_{\mathrm{I}}^{(i)},\boldsymbol{s}_{\mathrm{P}}^{(i)}$ by \eqref{eq:s_i}, \eqref{eq:s_p}
						\State Update $\boldsymbol{w}_{\mathrm{I/P},n}^{(i)}$, $\forall n$ by \eqref{M-eq:w}
						\State Compute $R^{(i)},z^{(i)}$ by \eqref{M-eq:R_waveform}, \eqref{M-eq:z_waveform}
					\Until $\lvert R^{(i)} - R^{(i-1)} \rvert \le \epsilon$ or $\lvert z^{(i)} - z^{(i-1)} \rvert \le \epsilon$
					\State Set $\boldsymbol{\phi}^{\star}=\boldsymbol{\phi}^{(i)},\boldsymbol{w}_{\mathrm{I/P}}^{\star}=\boldsymbol{w}_{\mathrm{I/P}}^{(i)}$
					\State \textbf{Output} $\boldsymbol{\phi}^{\star},\boldsymbol{w}_{\mathrm{I}}^{\star},\boldsymbol{w}_{\mathrm{P}}^{\star}$
				\end{algorithmic}
			\end{algorithm}
			Fig.~\ref{fi:re_noise}--\ref{fi:re_reflector} demonstrate that the LC-BCD Algorithm~\ref{al:lc_bcd} achieves near-optimal performance under different configurations.

			\begin{figure}[!h]
				\centering
				\subfloat[Average noise power\label{fi:re_noise}]{
					\resizebox{0.45\columnwidth}{!}{
						% This file was created by matlab2tikz.
%
%The latest updates can be retrieved from
%  http://www.mathworks.com/matlabcentral/fileexchange/22022-matlab2tikz-matlab2tikz
%where you can also make suggestions and rate matlab2tikz.
%
\definecolor{mycolor1}{rgb}{0.00000,0.44700,0.74100}%
\definecolor{mycolor2}{rgb}{0.85000,0.32500,0.09800}%
\definecolor{mycolor3}{rgb}{0.92900,0.69400,0.12500}%
\definecolor{mycolor4}{rgb}{0.49400,0.18400,0.55600}%
\definecolor{mycolor5}{rgb}{0.46600,0.67400,0.18800}%
%
\begin{tikzpicture}

\begin{axis}[%
width=4.521in,
height=1.575in,
at={(0.758in,2.472in)},
scale only axis,
xmin=0,
xlabel style={font=\color{white!15!black}},
xlabel={Average subband rate [bps/Hz]},
ymin=0,
ylabel style={font=\color{white!15!black}},
ylabel={Average output DC current [$\mu$A]},
axis background/.style={fill=white},
xmajorgrids,
ymajorgrids,
legend style={legend cell align=left, align=left, draw=white!15!black}
]
\addplot [color=mycolor1, line width=2.0pt, mark=o, mark options={solid, mycolor1}]
  table[row sep=crcr]{%
0.622682642041422	0\\
0.601236423383352	0.442612037820587\\
0.579761711004782	0.857115028670277\\
0.558290719126105	1.3002840100234\\
0.536819839004198	1.77053678956013\\
0.515349132207132	2.26589396248944\\
0.493878132430448	2.78475076674697\\
0.472407376006214	3.32546234426105\\
0.450936449697703	3.88676999650049\\
0.429464854645822	4.46766792652547\\
0.407993518307919	5.06691103182337\\
0.386521924519791	5.68388153522306\\
0.365049859056882	6.31804286614888\\
0.343577584389047	6.96894384519356\\
0.322105442370751	7.63606951807862\\
0.300632550078712	8.31998885518326\\
0.279159464190666	9.02026203791522\\
0.257685931948121	9.73753129549601\\
0.236212085740386	10.4394986349111\\
0.214738021476891	11.1977253766711\\
0.193263559348426	11.988742494465\\
0.171789496067895	12.83583619407\\
0.15031607724049	13.9277662691858\\
0.128843316789918	15.3451122069557\\
0.107370359882918	17.1694401639499\\
0.0858975730487439	19.5391294761626\\
0.0644244830789251	22.6441712726074\\
0.0429508607439586	26.9675178724178\\
0.0214765471232306	33.6400632006213\\
5.06778216560452e-10	55.2498608411414\\
};
\addlegendentry{$\sigma_n = -20$ dBm}

\addplot [color=mycolor2, dashed, line width=2.0pt, mark=+, mark options={solid, mycolor2}]
  table[row sep=crcr]{%
2.56470642681832	0\\
2.47626831055351	0.691282640556407\\
2.387830186357	1.44513790626809\\
2.29939207008782	2.2409281055306\\
2.21095402675715	3.06083609866881\\
2.12251603799647	3.89091474290932\\
2.03407822232999	4.72035471652295\\
1.94564039571933	5.5409962569854\\
1.85720290865869	6.34677356338329\\
1.76876562935097	7.13339325386499\\
1.6803284894462	7.89803703165199\\
1.59189208049175	8.63909403158659\\
1.50345552402019	9.3448319100033\\
1.4150199229818	10.0387783384305\\
1.32658411352241	10.7710496769656\\
1.23814884948203	11.6032048905254\\
1.14971343224388	12.512619603015\\
1.06127777598699	13.5921667321839\\
0.972841563324578	14.785292588885\\
0.884404361061965	16.1443219805427\\
0.795965783137808	17.648354099216\\
0.707527837245262	19.3174955630911\\
0.619092170622039	21.1418375993355\\
0.530657657556759	23.1962248898976\\
0.442221593215359	25.5356769719498\\
0.353786526568767	28.2145276766198\\
0.265350622047181	31.3837489123374\\
0.176909832311022	35.3145862426546\\
0.0884629483033547	40.7306098421577\\
7.09458767637855e-09	55.2514897021657\\
};
\addlegendentry{$\sigma_n = -30$ dBm}

\addplot [color=mycolor3, dotted, line width=2.0pt, mark=square, mark options={solid, mycolor3}]
  table[row sep=crcr]{%
5.60300411213724	0\\
5.40979707331538	1.26843811177322\\
5.21659003536169	2.65385979750085\\
5.02338299708889	4.0602207494935\\
4.83017595911997	5.42751271379642\\
4.63696892252802	6.72039590229198\\
4.44376189075816	7.92020071651595\\
4.25055486625106	8.9950008677074\\
4.05734786428981	9.98060775158619\\
3.86414089721892	10.9968594302406\\
3.67093400701335	12.1648985760136\\
3.47772719142908	13.548935954263\\
3.28452047603366	15.079092340491\\
3.09131397743405	16.7194872271433\\
2.89810758096324	18.4460754171478\\
2.70490118297073	20.2459759098262\\
2.51169514328775	22.09261857596\\
2.31848908462228	23.9780130283661\\
2.12528259502642	25.8914433044044\\
1.93207676715409	27.8250153090131\\
1.73887165719365	29.7731852303823\\
1.54566619427866	31.7395746582437\\
1.35246024801155	33.7311369746514\\
1.15925512492905	35.7581935450606\\
0.966051654123267	37.8390494500372\\
0.772850338062731	40.0063224185348\\
0.579652258714572	42.3203089394301\\
0.386458772722115	44.9028452278983\\
0.193262507567315	48.0740503031957\\
8.46957263577546e-08	55.2522464484787\\
};
\addlegendentry{$\sigma_n = -40$ dBm}

\addplot [color=mycolor4, dashdotted, line width=2.0pt, mark=x, mark options={solid, mycolor4}]
  table[row sep=crcr]{%
8.89182888246906	0\\
8.58521409326477	1.96098834911378\\
8.27859930431421	4.05989105460428\\
7.97198451540705	6.06870755780236\\
7.66536972726048	7.88362863064171\\
7.35875494234336	9.43618959493692\\
7.05214016485443	10.8127072170366\\
6.74552540774476	12.4741684254563\\
6.43891067407054	14.5555443991645\\
6.13229598553153	16.8454799083811\\
5.82568129342707	19.2731113712525\\
5.51906660032758	21.7521854553867\\
5.21245185016244	24.2323824942928\\
4.90583701789151	26.6734795248149\\
4.59922211455772	29.0440750052004\\
4.29260724906173	31.3225893835949\\
3.98599239438153	33.494767742988\\
3.67937755072139	35.552209873297\\
3.37276277265554	37.4914932566724\\
3.06614802067162	39.3132279975345\\
2.75953331578265	41.0214965274521\\
2.45291862026744	42.6229668131923\\
2.14630391889746	44.1271804350843\\
1.83968928821506	45.5461924021602\\
1.53307479104111	46.8954585545062\\
1.22646079951217	48.1952726796818\\
0.91984920608438	49.475365369516\\
0.613244610257377	50.7900023468271\\
0.306658752243681	52.2704828331503\\
6.6181262262034e-07	55.2524060040588\\
};
\addlegendentry{$\sigma_n = -50$ dBm}

\addplot [color=mycolor5, line width=2.0pt, mark=triangle, mark options={solid, mycolor5}]
  table[row sep=crcr]{%
12.2103843782899	0\\
11.7893366410274	2.68351453934089\\
11.3682889039	5.45483815607794\\
10.947241166918	7.92794126529428\\
10.5261934307476	9.95702816188572\\
10.1051456973549	11.8429216334637\\
9.68409797010892	14.5286151230504\\
9.26305024885062	17.622066442443\\
8.84200252584905	20.885258126308\\
8.4209547804507	24.1687916303299\\
7.99990702893788	27.3673003036257\\
7.57885928223634	30.41127222453\\
7.15781153964667	33.2562718591345\\
6.7367638062682	35.8772207913038\\
6.31571607994176	38.2640213719196\\
5.89466835672578	40.4174521525468\\
5.47362063871286	42.3459255694841\\
5.05257293876591	44.062920439037\\
4.63152521066552	45.5848791625207\\
4.21047753365302	46.929914855262\\
3.78942990876459	48.116631980624\\
3.36838239815027	49.1634333799863\\
2.9473350693401	50.0883433431305\\
2.52628811920357	50.9088251018236\\
2.10524175679114	51.6419592760158\\
1.68419753993712	52.3045265139519\\
1.26315258863362	52.9150424715196\\
0.842103490569631	53.5003838164046\\
0.421065549028916	54.1153815172928\\
2.6080731753108e-06	55.2525883066671\\
};
\addlegendentry{$\sigma_n = -60$ dBm}

\end{axis}

\begin{axis}[%
width=4.521in,
height=1.575in,
at={(0.758in,0.481in)},
scale only axis,
xmin=0,
xlabel style={font=\color{white!15!black}},
xlabel={Average subband rate [bps/Hz]},
ymin=0,
ylabel style={font=\color{white!15!black}},
ylabel={Average power splitting ratio},
axis background/.style={fill=white},
xmajorgrids,
ymajorgrids
]
\addplot [color=mycolor1, line width=2.0pt, mark=o, mark options={solid, mycolor1}, forget plot]
  table[row sep=crcr]{%
0.622682642041422	2.22044604925031e-16\\
0.601236423383352	0.0546236495114467\\
0.579761711004782	0.0964545009725418\\
0.558290719126105	0.137498191229909\\
0.536819839004198	0.177873331395111\\
0.515349132207132	0.217566407350811\\
0.493878132430448	0.25659248076223\\
0.472407376006214	0.294979558835803\\
0.450936449697703	0.332729136590715\\
0.429464854645822	0.369849855983362\\
0.407993518307919	0.406348130696068\\
0.386521924519791	0.442241902837711\\
0.365049859056882	0.477534030986953\\
0.343577584389047	0.512240810836996\\
0.322105442370751	0.546395015564245\\
0.300632550078712	0.579983984873656\\
0.279159464190666	0.613049615192229\\
0.257685931948121	0.645590971437943\\
0.236212085740386	0.676568007560787\\
0.214738021476891	0.706164872353165\\
0.193263559348426	0.731106269948913\\
0.171789496067895	0.748040698737134\\
0.15031607724049	0.758931022693324\\
0.128843316789918	0.767716079839159\\
0.107370359882918	0.777604274930107\\
0.0858975730487439	0.791535823793534\\
0.0644244830789251	0.808031330391159\\
0.0429508607439586	0.83543509777309\\
0.0214765471232306	0.877998850731802\\
5.06778216560452e-10	0.999985032883765\\
};
\addplot [color=mycolor2, dashed, line width=2.0pt, mark=+, mark options={solid, mycolor2}, forget plot]
  table[row sep=crcr]{%
2.56470642681832	2.22044604925031e-16\\
2.47626831055351	0.0725237133154718\\
2.387830186357	0.139989893381437\\
2.29939207008782	0.203190010398952\\
2.21095402675715	0.262397568599277\\
2.12251603799647	0.317864920758931\\
2.03407822232999	0.369827729912416\\
1.94564039571933	0.418502741913097\\
1.85720290865869	0.464099928347048\\
1.76876562935097	0.506809146273129\\
1.6803284894462	0.546807630682464\\
1.59189208049175	0.584283049025273\\
1.50345552402019	0.618110623271806\\
1.4150199229818	0.647394330204554\\
1.32658411352241	0.670893774727408\\
1.23814884948203	0.690553072971344\\
1.14971343224388	0.702701986186805\\
1.06127777598699	0.716148089469128\\
0.972841563324578	0.727078931470642\\
0.884404361061965	0.739903954550748\\
0.795965783137808	0.753663476334433\\
0.707527837245262	0.769233707717023\\
0.619092170622039	0.781201543791352\\
0.530657657556759	0.794418338288306\\
0.442221593215359	0.81249848992407\\
0.353786526568767	0.82999498237396\\
0.265350622047181	0.851744742483584\\
0.176909832311022	0.877782883033547\\
0.0884629483033547	0.913510037776107\\
7.09458767637855e-09	0.999982799897391\\
};
\addplot [color=mycolor3, dotted, line width=2.0pt, mark=square, mark options={solid, mycolor3}, forget plot]
  table[row sep=crcr]{%
5.60300411213724	2.22044604925031e-16\\
5.40979707331538	0.127991916926748\\
5.21659003536169	0.239416990927143\\
5.02338299708889	0.336413625068893\\
4.83017595911997	0.420830440782178\\
4.63696892252802	0.494271888532176\\
4.44376189075816	0.558135082661057\\
4.25055486625106	0.613040094373272\\
4.05734786428981	0.656891364326641\\
3.86414089721892	0.685246304153997\\
3.67093400701335	0.700430855912646\\
3.47772719142908	0.713644190958766\\
3.28452047603366	0.72516085317073\\
3.09131397743405	0.737903773114186\\
2.89810758096324	0.749181879664\\
2.70490118297073	0.762366637344574\\
2.51169514328775	0.775294507005966\\
2.31848908462228	0.789172542623828\\
2.12528259502642	0.803664565393353\\
1.93207676715409	0.818173885976562\\
1.73887165719365	0.83208041597836\\
1.54566619427866	0.84556073679751\\
1.35246024801155	0.859374206159318\\
1.15925512492905	0.873428150035787\\
0.966051654123267	0.887865763572983\\
0.772850338062731	0.902725298727963\\
0.579652258714572	0.918377613608431\\
0.386458772722115	0.935537485994346\\
0.193262507567315	0.956079644597371\\
8.46957263577546e-08	0.999981676713924\\
};
\addplot [color=mycolor4, dashdotted, line width=2.0pt, mark=x, mark options={solid, mycolor4}, forget plot]
  table[row sep=crcr]{%
8.89182888246906	2.22044604925031e-16\\
8.58521409326477	0.191487432238718\\
8.27859930431421	0.345503736471672\\
7.97198451540705	0.46925839785151\\
7.66536972726048	0.568558598749526\\
7.35875494234336	0.647110730127443\\
7.05214016485443	0.696472565086636\\
6.74552540774476	0.712655672424247\\
6.43891067407054	0.727540366310054\\
6.13229598553153	0.740785667038446\\
5.82568129342707	0.75766832428581\\
5.51906660032758	0.775827335215796\\
5.21245185016244	0.794042717894206\\
4.90583701789151	0.812077054968388\\
4.59922211455772	0.829266540690191\\
4.29260724906173	0.845381261669418\\
3.98599239438153	0.860389789296263\\
3.67937755072139	0.874325290693399\\
3.37276277265554	0.887229794231026\\
3.06614802067162	0.899178257492019\\
2.75953331578265	0.910246607268167\\
2.45291862026744	0.920528559060485\\
2.14630391889746	0.930117604171488\\
1.83968928821506	0.939112975882499\\
1.53307479104111	0.947644054145816\\
1.22646079951217	0.955854596069092\\
0.91984920608438	0.9639479491882\\
0.613244610257377	0.972263740645756\\
0.306658752243681	0.981600147939692\\
6.6181262262034e-07	0.999983100248849\\
};
\addplot [color=mycolor5, line width=2.0pt, mark=triangle, mark options={solid, mycolor5}, forget plot]
  table[row sep=crcr]{%
12.2103843782899	2.22044604925031e-16\\
11.7893366410274	0.252525958601193\\
11.3682889039	0.439875951856179\\
10.947241166918	0.578499985739617\\
10.5261934307476	0.678980367517115\\
10.1051456973549	0.713937979894303\\
9.68409797010892	0.729747106425759\\
9.26305024885062	0.748420055249861\\
8.84200252584905	0.772587079654286\\
8.4209547804507	0.797784052008363\\
7.99990702893788	0.821673256514771\\
7.57885928223634	0.843493890712623\\
7.15781153964667	0.863107409000432\\
6.7367638062682	0.880575320498125\\
6.31571607994176	0.896043652119298\\
5.89466835672578	0.909677777423865\\
5.47362063871286	0.921666673422102\\
5.05257293876591	0.93219104083251\\
4.63152521066552	0.941415475340555\\
4.21047753365302	0.949511297321253\\
3.78942990876459	0.95662186335818\\
3.36838239815027	0.962881087104149\\
2.9473350693401	0.968411662140905\\
2.52628811920357	0.973327538539347\\
2.10524175679114	0.977735249322931\\
1.68419753993712	0.98174608824415\\
1.26315258863362	0.985478138623143\\
0.842103490569631	0.989097850258867\\
0.421065549028916	0.992935851267291\\
2.6080731753108e-06	0.999988512921205\\
};
\end{axis}
\end{tikzpicture}%
					}
				}
				\subfloat[AP-IRS horizontal distance\label{fi:re_distance}]{
					\resizebox{0.45\columnwidth}{!}{
						% This file was created by matlab2tikz.
%
%The latest updates can be retrieved from
%  http://www.mathworks.com/matlabcentral/fileexchange/22022-matlab2tikz-matlab2tikz
%where you can also make suggestions and rate matlab2tikz.
%
\definecolor{mycolor1}{rgb}{0.00000,0.44700,0.74100}%
\definecolor{mycolor2}{rgb}{0.85000,0.32500,0.09800}%
\definecolor{mycolor3}{rgb}{0.92900,0.69400,0.12500}%
\definecolor{mycolor4}{rgb}{0.49400,0.18400,0.55600}%
\definecolor{mycolor5}{rgb}{0.46600,0.67400,0.18800}%
%
\begin{tikzpicture}

\begin{axis}[%
width=4.521in,
height=3.566in,
at={(0.758in,0.481in)},
scale only axis,
xmin=0,
xlabel style={font=\color{white!15!black}},
xlabel={Average subband rate [bps/Hz]},
ymin=0,
ylabel style={font=\color{white!15!black}},
ylabel={Average output DC current [$\mu$A]},
axis background/.style={fill=white},
xmajorgrids,
ymajorgrids,
legend style={legend cell align=left, align=left, draw=white!15!black}
]
\addplot [color=mycolor1, line width=2.0pt, mark=o, mark options={solid, mycolor1}]
  table[row sep=crcr]{%
5.62971756911577	0\\
5.43558937646659	1.33252799107006\\
5.24146118497288	2.7986795714249\\
5.04733299309185	4.29198893697339\\
4.85320480179634	5.74601077686804\\
4.65907661345474	7.12163050200202\\
4.46494842883311	8.39820581189421\\
4.27082026275745	9.53366125558458\\
4.07669211556099	10.5797051135205\\
3.88256401746303	11.6920496754887\\
3.68843602308717	12.9710500359586\\
3.49430816753776	14.4974171662899\\
3.30018042908507	16.1702388525605\\
3.10605317749813	17.9532763003293\\
2.91192589876244	19.8350654963931\\
2.7177987160973	21.7880322214965\\
2.52367139872176	23.792414204289\\
2.32954383228037	25.8340011729511\\
2.13541648786422	27.8990437784181\\
1.94128974567269	29.9821840417189\\
1.74716255671205	32.08278950124\\
1.55303566298019	34.2022429021717\\
1.35890896739747	36.3438523851885\\
1.16478225794241	38.5190086993999\\
0.970657441863005	40.7479905121655\\
0.776534800060461	43.0654408420295\\
0.582415639004176	45.5347624161656\\
0.38830114405591	48.2851042615184\\
0.194182852323423	51.6573336243653\\
7.79271316870932e-08	59.2703798102306\\
};
\addlegendentry{$d_H = 1$ m}

\addplot [color=mycolor2, dashed, line width=2.0pt, mark=+, mark options={solid, mycolor2}]
  table[row sep=crcr]{%
5.52514023798568	0\\
5.33461816034897	1.23273331799257\\
5.14409608359791	2.57955507507106\\
4.95357400662628	3.94707508500453\\
4.76305193040024	5.27693487101075\\
4.57252985880187	6.53477420320902\\
4.38200779461748	7.7024434052501\\
4.19148574595849	8.75020236521257\\
4.00096374912918	9.70605353577857\\
3.81044182631828	10.6929992652021\\
3.61992000280051	11.8386772871507\\
3.42939849960544	13.1908524285036\\
3.23887727721874	14.6844716541182\\
3.04835618887448	16.282932872661\\
2.8578356116359	17.9623550318536\\
2.66731453284533	19.7136155753527\\
2.47679377553042	21.5113251156866\\
2.28627256158812	23.3464744230169\\
2.09575078019503	25.208443517656\\
1.90523041718189	27.0893437087216\\
1.7147107635222	28.9847973160248\\
1.5241906844373	30.8971467824258\\
1.33366995321798	32.8332947409296\\
1.14315029717544	34.8045403173421\\
0.952632487343273	36.8277298186975\\
0.76211574717381	38.9348294930658\\
0.571601773590314	41.1837111803125\\
0.381093047411238	43.6922467546441\\
0.190580091996301	46.7713283860573\\
1.5855649339828e-06	53.7340774313462\\
};
\addlegendentry{$d_H = 2$ m}

\addplot [color=mycolor3, dotted, line width=2.0pt, mark=square, mark options={solid, mycolor3}]
  table[row sep=crcr]{%
5.37185480264753	0\\
5.18661842968689	1.10276860474339\\
5.00138205769196	2.29616523471063\\
4.81614568656564	3.50275595295885\\
4.63090931701202	4.67406840369621\\
4.44567295550072	5.78156253650663\\
4.26043662878479	6.81018036691245\\
4.07520035661757	7.74299873655263\\
3.88996414998658	8.58809996424829\\
3.70472821221378	9.43494618501664\\
3.5194926369243	10.407876944825\\
3.33425687697099	11.5404721289505\\
3.14902157804936	12.8133383113416\\
2.96378663175517	14.1789458245481\\
2.77855169555365	15.6116361147422\\
2.59331708679886	17.1060637764187\\
2.40808243923072	18.6481369478694\\
2.22284771264734	20.2235165695412\\
2.03761310866046	21.8254890576065\\
1.85237832748827	23.4472249187726\\
1.66714408480583	25.0873219820655\\
1.48191076791296	26.744648829895\\
1.29667621810219	28.4243702369402\\
1.11144285659041	30.1347323330858\\
0.926209727796795	31.8940083378838\\
0.740980293862373	33.7316008361788\\
0.555752953190297	35.6971716791412\\
0.370527991735988	37.8958571723873\\
0.185299369837713	40.6002442617824\\
2.23650569062746e-06	46.7264428411617\\
};
\addlegendentry{$d_H = 4$ m}

\addplot [color=mycolor4, dashdotted, line width=2.0pt, mark=x, mark options={solid, mycolor4}]
  table[row sep=crcr]{%
5.31502829750509	0\\
5.13175145947378	1.05901352989469\\
4.94847462213536	2.20127627133738\\
4.76519778727132	3.35444562759418\\
4.58192095498212	4.47320673968442\\
4.39864413809695	5.53089033208736\\
4.2153673877907	6.51341801702306\\
4.03209074337974	7.4151774895707\\
3.84881419109927	8.21362324059535\\
3.6655379814282	9.02061037991632\\
3.48226195349221	9.93751176683657\\
3.29898574146922	10.9969963517932\\
3.11571006540933	12.2005506919816\\
2.93243462766662	13.4902242808989\\
2.74915929577502	14.8452903585386\\
2.56588412673569	16.2529961583282\\
2.38260963853632	17.7126957770154\\
2.19933504265549	19.2038135614297\\
2.01605973895946	20.7211917089166\\
1.83278582061697	22.2575790318067\\
1.6495105058615	23.8121863387721\\
1.46623702749623	25.3843917773169\\
1.28296352674759	26.9798498541261\\
1.09968865770101	28.6055278796973\\
0.916415929992141	30.2769978175873\\
0.733145315789887	32.0261070169563\\
0.549876092927185	33.8987832442813\\
0.366611204391133	35.9941250164919\\
0.183341728257561	38.5736848231281\\
3.00580424622898e-06	44.4217569570318\\
};
\addlegendentry{$d_H = 6$ m}

\addplot [color=mycolor5, line width=2.0pt, mark=triangle, mark options={solid, mycolor5}]
  table[row sep=crcr]{%
5.28942153168888	0\\
5.1070276857406	1.04001905716027\\
4.92463384047887	2.16016645958468\\
4.74223999902017	3.29026764028289\\
4.55984617433321	4.38635316717864\\
4.37745234316496	5.42253933372905\\
4.19505862082739	6.38516994180292\\
4.01266505648352	7.26879907601144\\
3.83027155506989	8.04996478618108\\
3.64787838381793	8.84167636961934\\
3.4654852417826	9.73147830157504\\
3.2830920330011	10.7655868703336\\
3.10069940957405	11.9353836630271\\
2.91830720435323	13.1917314976326\\
2.73591478990386	14.5142328930052\\
2.55352320771064	15.8875641945599\\
2.37113140799457	17.3103101829804\\
2.18873988785427	18.7651015384578\\
2.00634828443065	20.2460468323242\\
1.82395692979393	21.7468749205771\\
1.64156529570313	23.2640332192424\\
1.45917350453234	24.8008519612018\\
1.27678351608541	26.3596084273071\\
1.09439223618738	27.9495772802957\\
0.91200116868915	29.5853397806375\\
0.729615013053546	31.2938014886304\\
0.54722975907755	33.125747479013\\
0.364847146807028	35.1762892874538\\
0.182459769165083	37.7021949556375\\
2.95411880011399e-06	43.4299683708002\\
};
\addlegendentry{$d_H = 7.5$ m}

\end{axis}
\end{tikzpicture}%
					}
				}
				\\
				\subfloat[Transmit antennas\label{fi:re_tx}]{
					\resizebox{0.45\columnwidth}{!}{
						% This file was created by matlab2tikz.
%
%The latest updates can be retrieved from
%  http://www.mathworks.com/matlabcentral/fileexchange/22022-matlab2tikz-matlab2tikz
%where you can also make suggestions and rate matlab2tikz.
%
\definecolor{mycolor1}{rgb}{0.00000,0.44700,0.74100}%
\definecolor{mycolor2}{rgb}{0.85000,0.32500,0.09800}%
\definecolor{mycolor3}{rgb}{0.92900,0.69400,0.12500}%
\definecolor{mycolor4}{rgb}{0.49400,0.18400,0.55600}%
\definecolor{mycolor5}{rgb}{0.46600,0.67400,0.18800}%
%
\begin{tikzpicture}

\begin{axis}[%
width=4.521in,
height=3.566in,
at={(0.758in,0.481in)},
scale only axis,
xmin=0,
xlabel style={font=\color{white!15!black}},
xlabel={Average subband rate [bps/Hz]},
ymin=0,
ylabel style={font=\color{white!15!black}},
ylabel={Average output DC current [$\mu$A]},
axis background/.style={fill=white},
xmajorgrids,
ymajorgrids,
legend style={legend cell align=left, align=left, draw=white!15!black}
]
\addplot [color=mycolor1, line width=2.0pt, mark=o, mark options={solid, mycolor1}]
  table[row sep=crcr]{%
5.50914838827946	0\\
5.31917775376218	1.21238415018016\\
5.12920712006203	2.53237434713813\\
4.93923648620775	3.87072056595873\\
4.74926585314339	5.17155623912808\\
4.55929522495152	6.4019337604629\\
4.36932460435472	7.54442965987761\\
4.17935399971781	8.58051451999257\\
3.98938344864228	9.49393206192153\\
3.79941297282513	10.4449974695456\\
3.60944259743222	11.5482026531884\\
3.41947256120648	12.8547904568156\\
3.22950281213667	14.2994951635005\\
3.03953319078188	15.8477053110534\\
2.84956411149923	17.4772022230921\\
2.65959455954133	19.1765527931851\\
2.46962522493195	20.9227426883177\\
2.27965547891626	22.7064304053862\\
2.08968519753875	24.5175977524265\\
1.89971635717942	26.3485393662365\\
1.70974811834305	28.1947945918833\\
1.51977953341583	30.0588513628518\\
1.32981021004462	31.9485579940201\\
1.13984192648845	33.8728277377253\\
0.949875845432248	35.8491157957908\\
0.759910653496104	37.9086273864258\\
0.569948171233342	40.1074108283297\\
0.379990602447721	42.5619962028736\\
0.190028583875184	45.578028320887\\
1.63891137764227e-06	52.4027633496848\\
};
\addlegendentry{$M = 1, L = 20$}

\addplot [color=mycolor2, dashed, line width=2.0pt, mark=+, mark options={solid, mycolor2}]
  table[row sep=crcr]{%
6.46189873526159	0\\
6.23907464042722	2.73723959929336\\
6.01625054658382	6.01987480905357\\
5.79342645229771	9.47936996496051\\
5.57060235843012	12.8897796145773\\
5.34777826472142	16.1201587119815\\
5.12495417340215	19.1019371373066\\
4.90213008936536	21.631149830557\\
4.67930600688948	24.3460462059138\\
4.45648195743752	27.7094305898286\\
4.23365790829412	31.8404343394148\\
4.01083391230115	36.4680743353983\\
3.7880099818419	41.418903535318\\
3.56518612806095	46.5912021403416\\
3.34236218293948	51.9392639192043\\
3.11953814885029	57.3912235139604\\
2.89671403345102	62.9025009809857\\
2.67388999187562	68.4355872838161\\
2.45106593830422	73.9614345462637\\
2.22824199991204	79.4619653726271\\
2.00541789538026	84.9273480334271\\
1.78259387154366	90.3578141323815\\
1.55977018263179	95.7652746152504\\
1.33694680303105	101.175808249317\\
1.11412389726573	106.636071206736\\
0.891301011824557	112.224900218251\\
0.668480992178671	118.086028912171\\
0.445665011213708	124.5081383894\\
0.222853670486381	132.25498170037\\
1.99460263980914e-08	149.492785651764\\
};
\addlegendentry{$M = 2, L = 20$}

\addplot [color=mycolor3, dotted, line width=2.0pt, mark=square, mark options={solid, mycolor3}]
  table[row sep=crcr]{%
7.35669951826879	0\\
7.10302022415528	6.20695838254019\\
6.84934093071593	14.5555150709397\\
6.59566163702678	23.6766693832912\\
6.34198234332767	32.7550836230748\\
6.08830304966753	41.3328722174662\\
5.83462375612751	49.1392393654032\\
5.58094446313422	55.8653798217785\\
5.32726517237947	65.6302151933533\\
5.07358588237067	78.1136870014957\\
4.81990659564493	92.1002689046529\\
4.56622729499044	107.116906156919\\
4.31254799927282	122.826848866696\\
4.05886870647144	138.996067378272\\
3.80518940974032	155.416860397166\\
3.55151011482067	171.915548506208\\
3.29783081976253	188.350381672684\\
3.04415152550868	204.610110942687\\
2.79047223159697	220.613072675702\\
2.53679293788656	236.306164338218\\
2.28311364442439	251.664248254292\\
2.02943435096141	266.690744843053\\
1.77575505752962	281.419721976425\\
1.52207576431207	295.922027696879\\
1.26839648013976	310.318563303724\\
1.01471727849096	324.808151289411\\
0.761038467157238	339.736675131913\\
0.507361550712717	355.792353558055\\
0.253689569734145	374.785607066888\\
2.10823254814809e-09	416.256686325004\\
};
\addlegendentry{$M = 4, L = 20$}

\addplot [color=mycolor4, dashdotted, line width=2.0pt, mark=x, mark options={solid, mycolor4}]
  table[row sep=crcr]{%
7.89718716435744	0\\
7.62487036522532	10.3237068848587\\
7.35255356675235	25.1926293047823\\
7.08023676800697	41.7190780171817\\
6.80791996925535	58.2148953649991\\
6.53560317052696	73.7527375673274\\
6.26328637185538	87.2714623973032\\
5.99096957371599	100.951262556214\\
5.71865277584552	122.571571344387\\
5.44633597819243	148.145893283286\\
5.17401917800472	176.057935773479\\
4.90170237819323	205.626472804135\\
4.62938557918985	236.307908454541\\
4.35706878040348	267.633070250063\\
4.08475198167879	299.198435611311\\
3.81243518291843	330.666820968081\\
3.54011838419245	361.766674372583\\
3.26780158546565	392.28977879866\\
2.99548478676102	422.087902710721\\
2.72316798817262	451.069257777827\\
2.45085118955515	479.195748238163\\
2.17853439092251	506.481916162961\\
1.90621759223693	532.997313212395\\
1.63390079351373	558.875513396727\\
1.36158399534812	584.333595825453\\
1.08926720608478	609.719850016771\\
0.81695051265075	635.624505674896\\
0.544634677967558	663.203322509861\\
0.272322224679745	695.477885815851\\
3.99906817076994e-10	765.047653590634\\
};
\addlegendentry{$M = 6, L = 20$}

\addplot [color=mycolor5, line width=2.0pt, mark=triangle, mark options={solid, mycolor5}]
  table[row sep=crcr]{%
8.36494421020624	0\\
8.07649785782375	17.0906432523215\\
7.78805150599832	43.4349226354144\\
7.49960515394812	73.1098964996926\\
7.21115880188307	102.729032880964\\
6.92271244984901	130.482089171242\\
6.63426609786475	153.103823907926\\
6.34581974609309	184.404455315327\\
6.05737339380685	228.95309474102\\
5.76892704253152	279.041577395493\\
5.48048068964625	333.030792139764\\
5.19203433741295	389.719782244971\\
4.90358798534091	448.061758998046\\
4.61514163329533	507.143846038953\\
4.32669528126863	566.192369447304\\
4.03824892921281	624.574052919183\\
3.74980257715848	681.791904363296\\
3.46135622514287	737.47719023877\\
3.17290987323309	791.379374309841\\
2.88446352124382	843.3559192679\\
2.59601716935408	893.363636162664\\
2.30757081731741	941.453504883204\\
2.01912446525775	987.771447076769\\
1.73067811318711	1032.56934839507\\
1.442231761333	1076.2366674772\\
1.15378541250716	1119.37248534768\\
0.8653390988339	1162.96529125834\\
0.576893004754171	1208.90845079472\\
0.288448898824029	1262.09156431317\\
1.32759421907693e-10	1375.29943262755\\
};
\addlegendentry{$M = 8, L = 20$}

\end{axis}
\end{tikzpicture}%
					}
				}
				\subfloat[Reflecting elements\label{fi:re_reflector}]{
					\resizebox{0.45\columnwidth}{!}{
						% This file was created by matlab2tikz.
%
%The latest updates can be retrieved from
%  http://www.mathworks.com/matlabcentral/fileexchange/22022-matlab2tikz-matlab2tikz
%where you can also make suggestions and rate matlab2tikz.
%
\definecolor{mycolor1}{rgb}{0.00000,0.44700,0.74100}%
\definecolor{mycolor2}{rgb}{0.85000,0.32500,0.09800}%
\definecolor{mycolor3}{rgb}{0.92900,0.69400,0.12500}%
\definecolor{mycolor4}{rgb}{0.49400,0.18400,0.55600}%
\definecolor{mycolor5}{rgb}{0.46600,0.67400,0.18800}%
%
\begin{tikzpicture}

\begin{axis}[%
width=4.521in,
height=3.566in,
at={(0.758in,0.481in)},
scale only axis,
xmin=0,
xlabel style={font=\color{white!15!black}},
xlabel={Average subband rate [bps/Hz]},
ymin=0,
ylabel style={font=\color{white!15!black}},
ylabel={Average output DC current [$\mu$A]},
axis background/.style={fill=white},
xmajorgrids,
ymajorgrids,
legend style={legend cell align=left, align=left, draw=white!15!black}
]
\addplot [color=mycolor1, line width=2.0pt, mark=o, mark options={solid, mycolor1}]
  table[row sep=crcr]{%
5.39398168528569	0\\
5.20798231612376	1.22068934125607\\
5.02198294841078	2.56950277527446\\
4.83598358031069	3.94504339005404\\
4.64998421548369	5.28407711465017\\
4.46398485764056	6.54942877072522\\
4.27798553126039	7.72144130681083\\
4.091986243543	8.7752876871929\\
3.90598709676167	9.74071738963898\\
3.71998810918807	10.7458731914163\\
3.53398949723431	12.010638947443\\
3.34799110781442	13.4346431937548\\
3.16199303634301	15.0037791719369\\
2.9759947374505	16.6753590900413\\
2.78999717951576	18.4253824654167\\
2.6039993003971	20.2429126752568\\
2.41800303993126	22.0993478965388\\
2.2320066213693	23.9888020825773\\
2.046009851812	25.9018778833084\\
1.86001408578827	27.8288584295265\\
1.67401901678915	29.7691546576733\\
1.48802158109351	31.7204258928399\\
1.30202381771407	33.6924372577175\\
1.11602625014286	35.6923496649217\\
0.930030812393232	37.7383792852281\\
0.74403768773903	39.8630695849754\\
0.558045288723006	42.125552446907\\
0.372053838720171	44.6406653795928\\
0.186057713724565	47.7219566310904\\
2.36932454981797e-06	54.7172968182554\\
};
\addlegendentry{$M = 1, L = 20$}

\addplot [color=mycolor2, dashed, line width=2.0pt, mark=+, mark options={solid, mycolor2}]
  table[row sep=crcr]{%
6.38939005284257	0\\
6.16906625750341	2.69287628984093\\
5.94874246292913	5.95082528791035\\
5.72841866804983	9.39307366824786\\
5.50809487323903	12.7873969053184\\
5.28777107891079	15.9999752865681\\
5.06744728549881	18.9611565604305\\
4.84712349547719	21.4864276452412\\
4.62679970994188	24.3316448730531\\
4.4064759398445	27.8656715446675\\
4.18615219636047	32.0493991155135\\
3.96582846588238	36.7434997655354\\
3.74550472804516	41.7559199866439\\
3.52518113267417	46.9774640185028\\
3.30485743916887	52.3629740854898\\
3.08453361835898	57.8492242308327\\
2.86420980675741	63.3856051619952\\
2.64388597983508	68.9353678592478\\
2.42356279142047	74.4703959123556\\
2.20323938505362	79.9696019972744\\
1.98291595605503	85.4231515850158\\
1.76259236068624	90.8355884421288\\
1.54226920137643	96.2170582818137\\
1.3219459417765	101.593889756663\\
1.10162335893682	107.012420776369\\
0.881301296260641	112.549536501417\\
0.660981736659531	118.348264205722\\
0.440666053512689	124.691980975248\\
0.220352992187177	132.330613664532\\
2.1203184642114e-08	149.312998721595\\
};
\addlegendentry{$M = 1, L = 40$}

\addplot [color=mycolor3, dotted, line width=2.0pt, mark=square, mark options={solid, mycolor3}]
  table[row sep=crcr]{%
7.43070367742958	0\\
7.17447251579894	6.76223494050246\\
6.91824135479948	16.0225139073115\\
6.66201019351539	26.1805852226025\\
6.40577903225118	36.2911655949891\\
6.14954787101543	45.8283591559144\\
5.89331670989148	54.2491141254435\\
5.6370855491387	62.6162727914728\\
5.38085439061751	74.0389608242132\\
5.12462323114612	88.4560971397017\\
4.86839207145889	104.518543166208\\
4.61216090795218	121.694160515101\\
4.35592974312383	139.613360953075\\
4.09969858088462	158.000731345434\\
3.84346741947112	176.622895040107\\
3.58723625822777	195.283720833368\\
3.33100509696229	213.823680843438\\
3.07477393576479	232.118479095191\\
2.8185427745383	250.078379785211\\
2.56231161339942	267.645727871055\\
2.30608045224395	284.795747297748\\
2.04984929120957	301.532866603133\\
1.793618129938	317.900066212108\\
1.53738697002282	333.975473915193\\
1.28115582337954	349.896043264659\\
1.02492478597244	365.879637297305\\
0.768694280513084	382.303920868734\\
0.512465816096588	399.921379429795\\
0.256241685289694	420.70263297094\\
1.01524736902585e-09	465.923450281919\\
};
\addlegendentry{$M = 1, L = 80$}

\addplot [color=mycolor4, dashdotted, line width=2.0pt, mark=x, mark options={solid, mycolor4}]
  table[row sep=crcr]{%
8.23880138607269	0\\
7.95470478628136	15.1034292760118\\
7.67060818695366	38.0861397821963\\
7.38651158746603	63.9035612748088\\
7.10241498798263	89.6638776069751\\
6.81831838848447	113.81579900219\\
6.53422178905932	133.897096443946\\
6.25012519006625	160.331974362531\\
5.96602859022063	198.465197550024\\
5.68193199043858	241.636182196109\\
5.39783539091209	288.216282398954\\
5.11373879138408	337.18320550881\\
4.82964219187496	387.63710929526\\
4.54554559239449	438.792723671441\\
4.26144899289152	489.983130047552\\
3.97735239340111	540.660956718383\\
3.69325579393235	590.395104529155\\
3.4091591944058	638.864257508494\\
3.12506259508036	685.84846477669\\
2.84096599575534	731.221682106316\\
2.55686939635633	774.943072674845\\
2.27277279689	817.054764246695\\
1.98867619749331	857.682676799196\\
1.70457959786	897.046588953194\\
1.42048299833069	935.487168236536\\
1.13638639946855	973.530871179723\\
0.852289819533533	1012.04747029419\\
0.56819359674422	1052.71426971506\\
0.284100445163735	1099.89249843865\\
5.63069237482087e-11	1200.52634620485\\
};
\addlegendentry{$M = 1, L = 120$}

\addplot [color=mycolor5, line width=2.0pt, mark=triangle, mark options={solid, mycolor5}]
  table[row sep=crcr]{%
8.84101363520959	0\\
8.53615109571498	28.7625878560105\\
8.23128855686155	75.8616032293243\\
7.92642601774056	129.38968578531\\
7.62156347861622	182.691348226463\\
7.31670093948751	232.31008339363\\
7.0118384004381	271.459476915312\\
6.70697586121548	343.07066409621\\
6.4021133220667	430.805069014299\\
6.09725078293508	527.760855838789\\
5.7923882437851	631.335275812944\\
5.48752570465246	739.225936188985\\
5.18266316551724	849.393692404821\\
4.87780062641785	960.083229275764\\
4.57293808732327	1069.83525672043\\
4.26807554822051	1177.48379802341\\
3.96321300908769	1282.14078304626\\
3.65835047004049	1383.1731858147\\
3.35348793116348	1480.1771810739\\
3.04862539208979	1572.95181834235\\
2.74376285297585	1661.47842505597\\
2.43890031382818	1745.90541473069\\
2.13403777468744	1826.53891097427\\
1.82917523555294	1903.87016503731\\
1.52431269648603	1978.60510517558\\
1.21945015748736	2051.78283734276\\
0.914587619880958	2125.06463581438\\
0.609725136883748	2201.54796518127\\
0.304864284502245	2289.18056397853\\
2.20539018341826e-11	2473.30891344661\\
};
\addlegendentry{$M = 1, L = 160$}

\end{axis}
\end{tikzpicture}%
					}
				}
				\caption{Average achievable R-E region under different configurations.}
			\end{figure}
			\label{re:2.7}
		\end{response}

		\begin{point}
			How practical is it to consider lossless reflection from the RIS? Specifically, by considering the magnitude to be 1, the reflection losses at the RIS have been ignored.
			\label{pt:2.8}
		\end{point}

		\begin{response}
			The assumption of lossless reflection is indeed ideal but common in microwave theory (e.g., \cite{Shen2020a}). Based on the circuit analysis, the authors of \cite{M-Abeywickrama2020} shows that the reflection amplitude depends on the phase shift, which is minimized at \num{0} phase shift and is negatively correlated to the element resistance. For $R_l=\SI{1}{\ohm}$, $\lvert{\phi_l}\rvert$ is \num{1} for $\theta_l=\pi$ but around \num{0.6} for $\theta_l=\num{0}$. From the perspective of field and propagation, \cite{Ozdogan2020} points out that if the dimension of IRS element is much larger than the wavelength, the beamwidth would become small and the scattered field strength would depend on the observation angle. However, to reduce the design complexity and provide a preliminary insight, we considered the lossless reflection pattern at the IRS in this paper.
			\label{re:2.8}
		\end{response}

		\begin{point}
			Minor comment: The size of all the numerical results figures is too small.
			\label{pt:2.9}
		\end{point}

		\begin{response}
			We sincerely apologize for the inconvenience. The figures have been amplified in the revised manuscript.
			\label{re:2.9}
		\end{response}
	\end{reviewer}

	\begin{reviewer}
		\begin{point}
			First of all, motivations of studying the IRS on SWIPT is very unclear to me. Please clarify.
			\label{pt:3.1}
		\end{point}

		\begin{response}
			We appreciate your suggestion and have modified the manuscript to emphasize this point. The major issue of WPT and SWIPT is the low harvested DC power. To increase it, boosting the harvester input power level would be very helpful. This is achieved by the IRS that can engineer the channels such that the power level at the receiver gets amplified. The smart channel control and low power consumption of IRS are expected to bring more opportunities to SWIPT. A major challenge for SWIPT is that information decoding and energy harvesting have different sensitivity on the received signal. Although conventional radios as Wi-Fi and Bluetooth can support a signal strength of \num{-90} to \SI{-100}{\dBm}, most existing harvesters can only capture signals at \num{-20} to \SI{-30}{\dBm} power level. Since the transmit power is strictly constrained by government regulations, it is crucial to increase the end-to-end power transfer efficiency to boost the harvested power and extend the system coverage. Therefore, we believe the passive beamforming ability of the IRS can significantly enhance WPT and SWIPT performance.
			\label{re:3.1}
		\end{response}

		\begin{point}
			Also, the contributions of this work are rather unclear, and thus, those should be better mentioned.
			\label{pt:3.2}
		\end{point}

		\begin{response}
			We have resummarized the contribution of this paper as follows.

			\emph{First}, we propose a novel IRS-aided SWIPT architecture based on joint waveform, active and passive beamforming design. To make an efficient use of the rectifier nonlinearity, we superpose a multi-carrier unmodulated power waveform (deterministic multisine) to a multi-carrier modulated information waveform and evaluate the performance under TS and PS receiving modes. The proposed joint waveform, active and passive beamforming architecture exploits the rectifier nonlinearity, a beamforming gain, and the channel selectivity across spatial and frequency domains to enlarge the achievable R-E region. This is the first paper to propose a joint waveform, active and passive beamforming architecture for IRS-aided SWIPT. By doing so, it is also the first paper to properly model the harvester nonlinearity (including its impact on both power and shape of the incoming waveform) and account for its crucial role in IRS-aided SWIPT.

			\emph{Second}, we characterize each R-E boundary point by energy maximization under rate constraint, and solve the problem by a Block Coordinate Descent (BCD) algorithm based on the Channel State Information at the Transmitter (CSIT). For active beamforming, we prove that the global optimal active information and power precoders coincide at Maximum-Ratio Transmission (MRT) even with rectifier nonlinearity. For passive beamforming, we propose a Successive Convex Approximation (SCA) algorithm and retrieve the IRS phase shift by eigen decomposition with optimality proof. Finally, the superposed waveform is optimized by the Geometric Programming (GP) technique. The IRS phase shift, active precoder, and waveform amplitude are updated iteratively until convergence. This is the first paper to jointly optimize waveform and active/passive beamforming in IRS-aided SWIPT.

			\emph{Third}, we introduce two closed-form adaptive waveform schemes to avoid the exponential complexity of the GP algorithm. The Water-Filling (WF) strategy for modulated waveform and the SMF strategy for multisine waveform are combined in the time and power domains to facilitate practical SWIPT implementation. To accommodate the suboptimal waveform schemes, we modify the passive beamforming algorithm and characterize the R-E region by varying the duration ratio under TS mode and the combining and splitting ratios under PS mode. The proposed low-complexity designs achieve near-optimal R-E performance in different scenarios.

			\emph{Fourth}, we provide numerical results to evaluate the proposed algorithms. It is concluded that 1) multisine waveform is beneficial to energy transfer especially when the number of subbands is large, 2) TS is preferred at low Signal-to-Noise Ratio (SNR) while PS is preferred at high SNR, 3) there exist two optimal IRS development locations, one close to the AP and one close to the UE, 4) the output SNR scales linearly with the number of transmit antennas and quadratically with the number of IRS elements, 5) due to the rectifier nonlinearity, the output DC current scales quadratically with the number of transmit antennas and quartically with the number of IRS elements, 6) for narrowband SWIPT, the optimal active and passive beamforming for any R-E point are also optimal for the whole R-E region, 7) for broadband SWIPT, the optimal active and passive beamforming depend on specific R-E tradeoff point and require adaptive designs, 8) the proposed algorithms are robust to practical impairments such as inaccurate cascaded CSIT and finite IRS reflection states.
			\label{re:3.2}
		\end{response}

		\begin{point}
			Please explain the derived results more intuitively for better understanding.
			\label{pt:3.3}
		\end{point}

		\begin{response}
			We have updated the discussion carefully and the reviewer is referred to the updated manuscript.
			\label{re:3.3}
		\end{response}

		\begin{point}
			Authors assumed the unrealistic situation: the channels are assumed to be perfectly known. However, in practice, the channel should be estimated, e.g., as studied in \cite{You2019,Kang2020}. It would be much better to discuss the channel estimation issue by citing the above references.
			\label{pt:3.4}
		\end{point}

		\begin{response}
			We have added those references and considered the influence of imperfect cascaded CSIT on the R-E performance. The reviewer is referred to Response \ref{re:1.1} for details. Simulation result in Fig.~\ref{fi:re_csi} demonstrates the robustness of the proposed joint waveform, active and passive beamforming algorithm to cascaded CSIT inaccuracy with different number of IRS elements.
			\label{re:3.4}
		\end{response}

		\begin{point}
			More simulation results should be added to better and aggregately validate the effectiveness of the proposed method.
			\label{pt:3.5}
		\end{point}

		\begin{response}
			Thank you for the suggestion. To address this comment, we have introduced two closed-form adaptive waveform design and compared the results with the original BCD Algorithm~\ref{M-al:bcd}. Fig.~\ref{fi:re_noise}--\ref{fi:re_reflector} demonstrate that the LC-BCD Algorithm~\ref{al:lc_bcd} achieves near-optimal performance under different configurations. Moreover, we have investigated in Fig.~\ref{fi:re_csi} and \ref{fi:re_quantization} the impact of imperfect cascaded CSIT and quantized IRS on the achievable R-E region. It can be concluded that the proposed algorithms are robust to CSIT inaccuracy of the cascaded link, and even $b=1$ (i.e. two-state reflection) brings considerable R-E gain over the benchmark scheme without IRS. These observations demonstrate the advantage of the proposed joint waveform, active and passive beamforming design in practical IRS-aided SWIPT systems.

			\begin{figure}[!h]
				\centering
				\resizebox{0.45\columnwidth}{!}{
					% This file was created by matlab2tikz.
%
%The latest updates can be retrieved from
%  http://www.mathworks.com/matlabcentral/fileexchange/22022-matlab2tikz-matlab2tikz
%where you can also make suggestions and rate matlab2tikz.
%
\definecolor{mycolor1}{rgb}{0.00000,0.44700,0.74100}%
\definecolor{mycolor2}{rgb}{0.85000,0.32500,0.09800}%
\definecolor{mycolor3}{rgb}{0.92900,0.69400,0.12500}%
\definecolor{mycolor4}{rgb}{0.49400,0.18400,0.55600}%
\definecolor{mycolor5}{rgb}{0.46600,0.67400,0.18800}%
\definecolor{mycolor6}{rgb}{0.30100,0.74500,0.93300}%
%
\begin{tikzpicture}

\begin{axis}[%
width=4.036in,
height=3.396in,
at={(0.677in,0.458in)},
scale only axis,
xmin=0,
xlabel style={font=\color{white!15!black}},
xlabel={Per-subband rate [bps/Hz]},
ymin=0,
ylabel style={font=\color{white!15!black}},
ylabel={Output DC current [$\mu$A]},
axis background/.style={fill=white},
xmajorgrids,
ymajorgrids,
legend style={legend cell align=left, align=left, draw=white!15!black},
title style={font=\huge}, label style={font=\huge}, ticklabel style={font=\LARGE}, legend style={font=\LARGE}
]
\addplot [color=mycolor1, line width=2.0pt, mark=o, mark options={solid, mycolor1}]
  table[row sep=crcr]{%
5.50989911497176	0\\
5.21990698812864	2.77802874336024\\
4.92991259001432	6.038869185882\\
4.63992083897796	9.23475278835535\\
4.3499306084325	12.1304078947137\\
4.05994122338663	14.6805637390731\\
3.76995619750306	17.8496080827396\\
3.47996789351475	21.8469544908054\\
3.1899844341134	26.3662353187812\\
2.89999986010502	31.1999780444902\\
2.61002374139079	36.2054910405585\\
2.32004269702453	41.2844251438971\\
2.03006919146576	46.3726809055504\\
1.74009087823982	51.4286563451092\\
1.45011191189263	56.4430685414929\\
1.1601361959374	61.4451144402281\\
0.870165622296932	66.5335679215567\\
0.580183943824222	71.910321436042\\
0.290197894466932	78.1304873928847\\
0.000129289748496336	91.1966119454087\\
};
\addlegendentry{No IRS}

\addplot [color=mycolor2, dashed, line width=2.0pt, mark=+, mark options={solid, mycolor2}]
  table[row sep=crcr]{%
6.19866473297976	2.88755997739766e-15\\
5.85746437828043	4.57837303631426\\
5.51678855455717	10.1618450361305\\
5.17651479161168	15.6324680255444\\
4.83693649776143	20.4424652164012\\
4.4982884656843	24.9923489865221\\
4.16035822772869	31.3813328021844\\
3.82327357839389	39.1316624095118\\
3.48728545623072	47.6854775648095\\
3.1524372842244	56.6607323308993\\
2.81911774170471	65.8011286956677\\
2.48835899896378	74.9374515937143\\
2.16040850072305	83.9380030221139\\
1.83588134140764	92.7321996641975\\
1.51497341396527	101.309660305128\\
1.19901846708282	109.722700833724\\
0.888808754667698	118.101225298078\\
0.584761800975399	126.796231366664\\
0.288265965065607	136.665117440505\\
1.67879226916162e-08	156.959655053007\\
};
\addlegendentry{$b = 1$}

\addplot [color=mycolor3, dotted, line width=2.0pt, mark=square, mark options={solid, mycolor3}]
  table[row sep=crcr]{%
6.43224082448431	3.27610369612353e-15\\
6.0898531645186	5.31212685976592\\
5.74753845020205	11.9206953634275\\
5.40541213618562	18.4486860175543\\
5.06336983964786	24.2233205759185\\
4.72152759738319	29.735175852791\\
4.37987865547185	37.4903786980729\\
4.03842965504395	46.9127351974931\\
3.69721922878753	57.3283609743573\\
3.35627538061948	68.2770323009597\\
3.01575850627015	79.4516172590205\\
2.67576803477419	90.6395330351928\\
2.33649173092397	101.67348451048\\
1.9980311149348	112.474981135374\\
1.66072372018151	123.02718233145\\
1.32475969898232	133.398621412332\\
0.99044694078195	143.76963486418\\
0.657922946777145	154.538410390468\\
0.32768284742819	166.780340746571\\
2.09156870371462e-08	192.025954075574\\
};
\addlegendentry{$b = 2$}

\addplot [color=mycolor4, dashdotted, line width=2.0pt, mark=x, mark options={solid, mycolor4}]
  table[row sep=crcr]{%
6.50707708244435	3.41852479659429e-15\\
6.16435082250026	5.59024447407167\\
5.82163259683868	12.5939118843714\\
5.47891806602544	19.5306673271929\\
5.13621422113226	25.6788213638393\\
4.79351954702848	31.565834995333\\
4.45083763369798	39.8571761024993\\
4.10817557816947	49.9341292385325\\
3.76553428618393	61.0827330086535\\
3.42291382384105	72.8079760849851\\
3.08030452225355	84.7845801523822\\
2.73775412665916	96.7769668037776\\
2.39523959555037	108.61385784728\\
2.05275088132373	120.207391898607\\
1.71033849624943	131.545610877493\\
1.36800892217148	142.701587993965\\
1.02580820237254	153.857192341203\\
0.683731887881355	165.440160780643\\
0.341802268625104	178.623385444774\\
2.24911409862557e-08	205.84064483428\\
};
\addlegendentry{$b = 4$}

\addplot [color=mycolor5, line width=2.0pt, mark=triangle, mark options={solid, mycolor5}]
  table[row sep=crcr]{%
6.51219310805307	3.42854431433599e-15\\
6.16944509167559	5.61009443843645\\
5.82669713663092	12.6420930033537\\
5.48394904095408	19.6080649686877\\
5.14120105227918	25.7831270676674\\
4.79845321404746	31.6979007222729\\
4.45570547926302	40.0278744817702\\
4.11295764195447	50.1522587431315\\
3.77021011801552	61.3531074190372\\
3.42746306818565	73.135060051636\\
3.084714893668	85.1693687702944\\
2.74196811914805	97.2204655091078\\
2.39922232147663	109.115889937508\\
2.05647512100103	120.766506819247\\
1.71373023166313	132.160158384373\\
1.37098481930642	143.368169419647\\
1.02824283022912	154.580318957371\\
0.685506941267945	166.223151371196\\
0.342777878482553	179.473827400745\\
2.26007858704721e-08	206.833019634685\\
};
\addlegendentry{$b = 8$}

\addplot [color=mycolor6, dashed, line width=2.0pt, mark=triangle, mark options={solid, rotate=180, mycolor6}]
  table[row sep=crcr]{%
6.51221287559059	2.22044604925031e-10\\
6.16946482931077	5.61017038726914\\
5.82671678342275	12.6422757618082\\
5.48396873815337	19.6083546697317\\
5.14122069624897	25.7835332901979\\
4.79847266633801	31.6983722953733\\
4.45572466929463	40.0284890543344\\
4.11297675770656	50.152981594314\\
3.77022883276422	61.3540477957114\\
3.42748101702145	73.1361545866818\\
3.08473323815122	85.1706918384703\\
2.7419855457683	97.2216283461067\\
2.39923738578064	109.117492933191\\
2.05648986346861	120.768893232785\\
1.71374324528809	132.162682513428\\
1.37099661863584	143.37091238812\\
1.02825259887354	154.583327186729\\
0.685513439808935	166.226325848819\\
0.342780008699142	179.477251618881\\
2.26023192316464e-08	206.837086736742\\
};
\addlegendentry{Unquantized IRS}

\end{axis}
\end{tikzpicture}%
				}
				\caption{Average R-E region with quantized IRS for $M=1$, $N=16$, $L=20$, $\sigma_n^2=\SI{-40}{\dBm}$, $B=\SI{10}{\MHz}$ and $d_{\mathrm{H}}=d_{\mathrm{V}}=\SI{2}{\meter}$.}
				\label{fi:re_quantization}
			\end{figure}
			\label{re:3.5}
		\end{response}

		\begin{point}
			The sizes of figures are too small.
			\label{pt:3.6}
		\end{point}

		\begin{response}
			We sincerely apologize for the inconvenience. The figures have been amplified in the revised manuscript.
			\label{re:3.6}
		\end{response}
	\end{reviewer}


	\bibliographystyle{IEEEtran}
	\bibliography{IEEEabrv,library.bib}
\end{document}
